
%!TEX TS-program = xelatex
%!TEX encoding = UTF-8 Unicode

\documentclass[english,a4paper,11pt]{article}
\usepackage[T1]{fontenc}
\usepackage[utf8]{inputenc}
\usepackage{}
\usepackage[margin=3cm,a4paper]{geometry}

\usepackage{graphicx}
\usepackage{amssymb}
\usepackage[norsk]{babel}
\usepackage{color}
\usepackage{xcolor}
\usepackage{booktabs}
\usepackage{array}
\usepackage{wrapfig}
\usepackage{calc}
\usepackage{hyperref}
\usepackage{xspace}
\usepackage{parskip}


\definecolor{xalblue}{rgb}{0.1992,0.1992,0.5977}
\definecolor{xaldark}{rgb}{0.1992,0.1992,0}
\definecolor{lightgray}{gray}{0.8}


%\usepackage{fontspec,xltxtra,xunicode,relsize}
%\defaultfontfeatures{Mapping=tex-text}
%\setmainfont[Mapping=tex-text]{Verlag Light}
%\setmainfont[Mapping=tex-text]{Sina-Regular}
%\setsansfont[Scale=MatchLowercase,Mapping=tex-text]{Helvetica Neue}
%\setmonofont[Scale=MatchLowercase]{Monaco}
\renewcommand*{\familydefault}{\sfdefault}

\newcolumntype{L}{>{\raggedright}p{0.16\textwidth}}
\newcolumntype{R}{p{0.74\textwidth}}
\newcommand\VRule{\color{lightgray}\vrule width 0.5pt}


\usepackage{titlesec}
\newcommand\secformat[1]{\normalfont\Large\bfseries\color{xalblue}#1}
\titleformat{name=\section}[block]{}{}{}{\secformat}

\begin{document}

\vspace*{-3cm}
\hspace*{12.5cm}\includegraphics[width=4cm]{logo.pdf}

\vspace*{-1.65cm}

\noindent
{\LARGE \bfseries \color{xaldark}Curriculum Vitae for Åsmund Ødegård}

\section*{Personal information}

\begin{minipage}[t]{0.4\textwidth}
\begin{tabbing}Nationality---\=olaf\kill 
Address:\>Sneisveien 26, \\
\>1407 Vinterbro \\
Born:\>23.08.1972 \\
\end{tabbing}
\end{minipage}
$\quad\quad$
\begin{minipage}[t]{0.3\textwidth}
\begin{tabbing}Nationality---\=olaf\kill 
E-mail: \>asmund@xal.no \\
Phone: \> +47 900 69 915 \\
Nationality:\>Norwegian\\
\end{tabbing}
\end{minipage}

\section*{Summary}
I am a senior software developer, project leader, and researcher, with 19 years
of professional experience.
Most of my
career prior to founding Expert Analytics AS together with Ola Skavhaug,
has been with Simula Research Laboratory,
ranked as Norway's best ICT research institution in 2012. Through various
positions in that company I have obtained a broad set of technical and management
skills, which I since 2013 has used as Managing Director and Senior Consultant
in Expert Analytics. My main areas of technical expertise spans from mathematical
software development and scripting, to administration of advanced ICT
solutions and configuration management. Both through education and
practical hands-on experience I have developed strong analytical skills within
my fields of expertise.
\\
\\
As manager and leader of several projects with internal and external
participants, I have learned the importance of understanding the strengths and
skills of the participants in order to deliver on time. Communication within
the team is always key to ensure that everyone have ownership in the project and
stay motivated. This must always be balanced with a firm leadership.


\section*{Technical Skills}
\begin{tabular}{LR}
Languages & Python, C, C++, C\#, PHP, Java, Javascript, Bash, Matlab, \LaTeX, HTML, Sql, CSS, Scala, Lisp \\
Frameworks & Google App Engine, Puppet, Numpy, Pandas, SciPy, MPI, BSP, Diffpack, Plone, Knockout, jQuery, RequireJS, Web.py, Turbogears, Bootstrap, Slim, Twig\\
Tools & Linux, OS X, VirtualBox, Vagrant, MySql, Sqlite, NoSqls, Apache, Nginx, Git, Subversion, CVS, Mercurial, Make, SCons, Autoconf, buildout
\end{tabular}

\section*{Education}
\begin{tabular}{LR}
2006 & Dr. Scient in Scientific Computing, Faculty of Mathematics and Natural Sciences, University of Oslo. Thesis' title ``Applications of high level software for parallel solution of Partial Differential Equations.'' \\
1995 & Cand. Scient in Applied and Industrial Mathematics, Department of Mathematics, University of Oslo
\end{tabular}

\section*{Professional Experience}
\begin{tabular}{LR}
2013 --      & Managing Director and Consultant at Expert Analytics \\
2011 -- 2013 & Director of External Research Funding at Simula Research Laboratory \\
2010 -- 2011 & Director of Administration at Simula Research Laboratory \\
2008 -- 2010 & Assistant Director of Administration at Simula Research Laboratory \\
2005 -- 2007 & Research Scientist at Simula Research Laboratory \\
2004 -- 2009 & IT Manager at Simula Research Laboratory \\
2001 -- 2005 & Scientific Programmer at Simula Research Laboratory \\
2000 -- 2001 & Scientific Programmer at Department of Informatics, University of Oslo \\
1996 -- 1999 & Research Assistant at Department of Informatics, University of Oslo \\
1996         & Summer intern, SINTEF Applied Mathematics \\
1993 -- 1994 & Assistant teacher, Department of Mathematics, University of Oslo \\
\end{tabular}

\section*{Languages}
\begin{tabular}{LR}
Norwegian & Mother tongue \\
English & Fluent \\
%German & Basic \\
%French & Rudimentary \\
\end{tabular}

\section*{Personal Skills}
\begin{tabular}{LR}
Management & Development and implementation of strategies, guidelines and procedures, lead and encourage staff, personnel follow-up \\
Communication & Writing and editing, typesetting \& typography, try to learn French \\
Funding & EU and RCN funding instruments, consortium development, proposal writing \\
Technology & Extensive knowledge about system administration; always eager to learn about new frameworks and open source technologies; solution architect and problem solver \\
\end{tabular}

\section*{Some interests and hobbies}
\begin{tabular}{LR}
Soft & Playing clarinet and saxophone, cooking \& wine, photography \\
Hard & Bicycle, motorcycle \\
Tech & Fonts, typesetting with \TeX, computers, gadgets, coding
\end{tabular}

\section*{Extended descriptions of select projects}

\begin{tabular}{@{}LR}
Activity & Demographic and financial estimations \\
Role & Architect and main developer \\
Staffing & Project leader, developer, front-end designer \\
Description & For customer we developed a web-based tool for statistical estimation of demographics within a large closed group, and a tool for economical planning of a large scale development. My main contribution was the mathematical analysis of the problem, and development of the necessary logic and algorithms. Further, I also implemented the system in a custom application server. The work was carried out in close collaboration with financial experts and a designer. \\
Tools & Python, Pandas, Web.py, Sqlite, MS Sql, Excel, HTML, CSS, JavaScript, jQuery  \\ 
\addlinespace \bottomrule[.1pt] \addlinespace
\end{tabular}

\begin{tabular}{@{}LR}
Activity & Establishing technical infrastructure of Simula Research Laboratory \\
Role & Project leader and main architect \\
Staffing & Project leader, 4 system administrators, external contractors \\
Description & For the newly created research institute we designed and developed a full technical infrastructure in the renovated office building at the old airport at Fornebu; office networks, back-end facilities and computer room, telephone systems, printers, servers, etc. Technical requirements, budgetary limitations and practical issues had to be taken into consideration in order to establish a scalable and flexible infrastructure.  \\ 
\addlinespace \bottomrule[.1pt] \addlinespace
\end{tabular}

\begin{tabular}{@{}LR}
Activity & Symphonical \\
Role & Main developer \\
Staffing & Project leader and two developers \\
Description & Symphonical is a web-based collaboration tool. Initially it was conceived as a tool for running agile software development projects based on the metaphor of post-it notes on a virtual wall, a scope that since has been widened. We created the first prototype of the system, in 2005, before it was spun out as a company and developed further by others. \\
Tools & PHP, MySql \\ 
\addlinespace \bottomrule[.1pt] \addlinespace
\end{tabular}

\begin{tabular}{@{}LR}
Activity & Simulation of Medical Ultrasound \\
Role & Researcher and main developer \\
Staffing & 4--5 researchers and developers \\
Description & A research project where we studied different mathematical models for numerical simulation of medical ultrasound. My main responsibility was implementation of object oriented finite element based simulators for ultrasound propagating in inhomogeneous media. The target application was detection of tumors in tissue. \\
Tools & C++, Diffpack, MPI, VTK \\ 
\addlinespace \bottomrule[.1pt] \addlinespace
\end{tabular}

\begin{tabular}{@{}LR}
Activity & PySE \\
Role & Main developer and researcher \\
Description & Python Stencil Environment (PySE) is a high level environment for solving partial differential equations using finite difference methods (FDM). The environment provides an abstraction for the stencil concept of FDM, which makes it possible to implement FDM-based simulators in notation very close to the mathematical formulation. The environment is inherently parallell and distributes a problem on the available hardware in an optimal way based on the communication patterns defined by the stencils. PySE was developed as part of my PhD studies and published as open source software.\\
Tools & Python, Numpy, Scipy, SWIG, C++, MPI \\ 
\addlinespace \bottomrule[.1pt] \addlinespace
\end{tabular}

\begin{tabular}{@{}LR}
Activity & Python Computing Components (PyCC) \\
Role & Developer \\
Description & PyCC is a modern and efficient scripting framework that is used to solve differential equations to study the electrical activity in the human heart. The main parts of PyCC is a high level scripting interface implemented in Python for flexibility and low-level extension modules implemented in C/C++ and Fortran for efficiency. My main role in the project was to provide a flexible build system for PyCC taking into account all the dependencies of the system. The build system was implemented using the SCons framework. \\
Tools & C/C++, Fortran, Python, Swig, MPI, Subversion, SCons, PETSc, FEniCS, Hypre, BoomerAMG, Diffpack, GNU Compiler Collection and Debugger  \\ 
\addlinespace \bottomrule[.1pt] \addlinespace
\end{tabular}

\begin{tabular}{@{}LR}
Activity & Web site for Simula Research Laboratory \\
Role & Main developer and designer \\
Description & In order to facilitates new communication strategies for the company, it was necessary to update the web site for Simula Research Laboratory. New channels for communicating news and ongoing research projects were created, as well as a complete redesign of the company Intranet which at the same time replaced mailing list as the primary form of official internal communication. The process of upgrading the system consisted both of technical work and meetings with all the stakeholders within the organisation in order to develop a system in line with their needs.\\
Tools & Python, PHP, Javascript, Plone, MySql \\ 
\addlinespace \bottomrule[.1pt] \addlinespace
\end{tabular}

\begin{tabular}{@{}LR}
Activity & Outsourcing of ICT services \\
Role & Project coordinator \\
Description & Simula Research Laboratory decided in 2012 to put a contract for its IT operations out for tender, after running the services with internal resources for a decade. My core responsibility was authoring the specification of the tender and evaluating the offers as well as participate in negotiating the final contract. I was also responsible for handling the communication with all the internal stakeholders. As the change affected the daily work of all employees it was particularly important to have an open and inclusive dialog in the process to ensure a satisfactory result. \\ 
\addlinespace \bottomrule[.1pt] \addlinespace
\end{tabular}

\begin{tabular}{@{}LR}
Activity & mCASH development\\
Role & Senior Python Developer \\
Staffing & 12--15 developers \\
Description & My responsibilities in this project included integrations of third-party APIs for user-identification with Norwegian BankID, and development of APIs for access to mCASH resources for external integrators, based on OpenID Connect and Oauth protocols. I was also deeply involved in design and implementation of the transactional engine handling the payments; the core of the mCASH system. This engine connects persons, merchants and banks together, and computes transation fees for involved parties in addition to moving correct amounts. The backend system of mCASH is implemented solely in Python. The development at mCASH followed the Scrum agile method.\\
Tools & Python, Google app engine, webapp2, Jinja2, Javascript, jQuery, HTML, Nose tests, Git, Jira. \\ 
\addlinespace \bottomrule[.1pt] \addlinespace
\end{tabular}

\begin{tabular}{@{}LR}
Activity & SmartBlikk development \\
Role & Systems Architect and Senior Developer \\
Description & SmartBlikk is a business intelligence solution for the SME market. My responsibility was development of an API for the system for import of data from external systems, as well as integrating external APIs for use from SmartBlikk, mainly for importing accounting data. The pilot of SmartBlikk was created by the founder of Intelliview, and I am in charge of designing and implementing the second version of the system with dynamic in-browser reporting capabilities. I am also in charge of establishing a deployment environment for SmartBlikk on the cloud, specifically on Amazon Elastic Beanstalk scaling out with multiple EC2 nodes and using AWS RDS database. \\
Tools & Php, Mysql, Slim, Twig, Javascript, jQuery, AWS EB, AWS EC2, AWS RDS. \\ 
\addlinespace \bottomrule[.1pt] \addlinespace
\end{tabular}

\begin{tabular}{@{}LR}
Activity & FrontSystems development \\
Role & Senior Developer and adviser \\
Staffing & 3 developers \\
Description & FrontSystems is a Norwegian company developing a web-based POS system. The current system was build on some technologies that are approaching EOL, and my main role was to advise and contribute towards a more modern architecture. This include moving development and deployment to Github and Azure as well as implementing reporting services in POS backend using client-side JS rather than Microsoft reporting services. \\
Tools & C\#, JavaScript, Asp.Net, Knockout, Bootstrap, RequireJS, jQuery, MSSQL, TeamCity, Git, Azure. \\ 
\addlinespace \bottomrule[.1pt] \addlinespace
\end{tabular}

\begin{tabular}{@{}LR}
Activity & Building linux clusters \\
Role & Project manager, lead architect and developer \\
Staffing & 2 additional Systems administrators \\
Description & For the scientific computing group at Ifi, University of Oslo and later at Simula Research Laboratory, I was in charge of building three different compute clusteres for advanced numerical computations. The systems grew from 32 compute cores up to 576, and was configured with core os, multiple networks, schedulers and queue systems. The first clusters where managed with CFengine as provisioning system, while the last one was provisioned using Puppet. \\
Tools & Debian GNU/Linux, CFengine, Puppet, OpenPBS, MAUI, MPIch, BSP, \\
\end{tabular}

\end{document}

